% SOTA Comparison Table for distantly supervised NER
% To be inserted in Section 4 or Discussion
\begin{table}[t]
\centering
\small
\begin{tabular}{lccc}
\toprule
\textbf{Method} & \textbf{CoNLL-2003} & \textbf{Webpage} & \textbf{Wikigold} \\
\midrule
\multicolumn{4}{c}{\textit{Supervised upper bound (RoBERTa)}} \\
\quad Fully Supervised (BOND) & 90.11 & 72.39 & 86.43 \\
\midrule
\multicolumn{4}{c}{\textit{Distant supervision + denoising methods}} \\
\quad BOND Stage I \cite{liang2020bond} & 75.61 & 59.11 & 52.15 \\
\quad BOND (full) \cite{liang2020bond} & 81.48 & 65.74 & 60.07 \\
\quad RoSTER \cite{meng2021roster} & 85.4 & -- & 67.8 \\
\quad SANTA \cite{shao2023santa} & 85.11 & 71.79 & -- \\
\quad ATSEN \cite{zhang2021atsen} & 85.59 & 70.55 & -- \\
\quad SCDL \cite{zhang2022scdl} & 83.69 & 68.47 & 64.13 \\
\midrule
\multicolumn{4}{c}{\textit{Distant supervision + hyperparameter scheduling (ours)}} \\
\quad Static baseline (roberta-base) & [TBD] & 61.5 & 54.2 \\
\quad GP-TS tuned (ours) & [TBD] & 62.7 & [TBD] \\
\quad TD3 tuned (ours) & -- & -- & 54.0 \\
\bottomrule
\end{tabular}
\caption{Comparison with state-of-the-art distantly supervised NER methods. All scores are entity-level micro-averaged F1 (\%). ``--'' indicates the method did not report results on that dataset. Our scheduling methods operate on the BOND Stage I training loop (roberta-base) without self-training or denoising, making them complementary to the denoising approaches above. [TBD] values will be filled upon experiment completion.}
\label{tab:sota-comparison}
\end{table}
